\textbf{Proof.}
Put \(N=2n\), index the signed exposure coordinates by \(j\), write \(\pi_j=\mathbb E_\nu D_j\), and put \(Z_j=D_j-\pi_j\).  Exposure positivity gives \(\pi_{\min,n}>0\), so every inverse probability below is well defined.  Known exposure gives \(Y_i^{\mathrm{obs}}=Y_i(a)\) on \(D_{i,a}=1\).  With the signed convention in \(Y\), this yields the exact identities
\[
\mathbb E_\nu\widehat Y_j^{\mathrm{obs}}=Y_j,
\qquad
\{\Pi^{-1}(\widehat Y^{\mathrm{obs}}-Y)\}_j
=\frac{Y_jZ_j}{\pi_j^2}.
\tag{1}
\]

If \(\operatorname{rank}(S_B)=0\), then the positive-semidefinite matrix \(S_B\) is zero and its spectral pseudoinverse is zero.  Consequently
\[
\widehat\gamma^{\mathrm{HT}}(B)=0=\gamma_B^{*,+}(Y).
\tag{2}
\]
Substitution into the two displayed estimator maps gives \(n^{-1}\mathbf1^{\top}\widehat Y^{\mathrm{obs}}\) in both cases.  This proves every rank-zero assertion, including the absence of control-variate gain.

Now fix \(r\in\{1,2\}\) and \([B]_+\in\mathcal S_{r,n}(\lambda_r,\overline a)\).  Put
\[
G_B=\widehat\gamma^{\mathrm{HT}}(B)-\gamma_B^{*,+}(Y).
\]
By (1), with one consistently named coefficient vector,
\[
G_B=\sum_{j=1}^{N}u_{B,j}Z_j,
\qquad
u_{B,j}=S_B^+B^{\top}\Omega e_j\frac{Y_j}{\pi_j^2}.
\tag{3}
\]
The positive-eigenvalue margin gives
\[
\lVert S_B^+\rVert_{\mathrm{op}}\le(n\lambda_r)^{-1}.
\tag{4}
\]
Every column of \(\Omega\) has at most \(s_n\) nonzero entries, every covariance entry has absolute value at most one, the outcomes have absolute value at most one, and the normalized basis coordinates are bounded by \(\overline a\).  Hence
\[
\lVert B^{\top}\Omega e_j\rVert_2
\le\sqrt2\,\overline a s_n,
\qquad
\sup_j\lVert u_{B,j}\rVert_2\le U_{n,r}.
\tag{5}
\]
The centered basis statistic is
\[
X_B=n^{-1}B^{\top}(D-\mathbb E_\nu D)
=\sum_{j=1}^{N}x_{B,j}Z_j,
\qquad
x_{B,j}=n^{-1}B^{\top}e_j,
\tag{6}
\]
so \(\sup_j\lVert x_{B,j}\rVert_2\le X_n\).  Applying the fourth-moment inequality in \(\texttt{def:rank-adaptive-stability-strata}\) to each of the two coordinates, and using \((u^2+v^2)^2\le2(u^4+v^4)\), gives
\[
\mathbb E_\nu\lVert G_B\rVert_2^4
\le4\mathfrak M_{4,n}U_{n,r}^4,
\qquad
\mathbb E_\nu\lVert X_B\rVert_2^4
\le4\mathfrak M_{4,n}X_n^4.
\tag{7}
\]
The feasible and oracle estimators satisfy
\[
\widehat\tau_B^{\mathrm{CV}}
=\widehat\tau_B^{\mathrm{or},+}+\Delta_B,
\qquad
\Delta_B=-G_B^{\top}X_B.
\tag{8}
\]
Cauchy--Schwarz and (7) imply
\[
\mathbb E_\nu\Delta_B^2
\le4\mathfrak M_{4,n}U_{n,r}^2X_n^2.
\tag{9}
\]

Put \(T_B=\widehat\tau_B^{\mathrm{or},+}-\tau\).  Equation (1) makes the Horvitz--Thompson term unbiased, and the basis statistic is centered by construction, so \(\mathbb E_\nu T_B=0\).  Direct expansion gives
\[
T_B=\sum_{j=1}^{N}t_{B,j}Z_j,
\qquad
t_{B,j}=n^{-1}\left\{\frac{Y_j}{\pi_j}
-B_{j\cdot}\gamma_B^{*,+}(Y)\right\}.
\tag{10}
\]
Write \(v=\Omega^{1/2}\Pi^{-1}Y\) and \(Z_B^\circ=\Omega^{1/2}B\).  The projector identity in \(\texttt{lem:rank-stratified-projector-reduction}\) gives
\[
\gamma_B^{*,+}(Y)^{\top}S_B\gamma_B^{*,+}(Y)
\le Y^{\top}\Pi^{-1}\Omega\Pi^{-1}Y.
\tag{11}
\]
There are at most \(2ns_n\) nonzero entries in the covariance sum, so bounded outcomes and positivity give
\[
Y^{\top}\Pi^{-1}\Omega\Pi^{-1}Y
\le2ns_n\pi_{\min,n}^{-2}.
\tag{12}
\]
The vector \(\gamma_B^{*,+}(Y)\) belongs to the positive-eigenvalue range of \(S_B\).  Combining (11), (12), and \(S_B\succeq n\lambda_r\) on that range yields
\[
\lVert\gamma_B^{*,+}(Y)\rVert_2\le\Gamma_{n,r},
\qquad
\sup_j|t_{B,j}|\le T_{n,r}.
\tag{13}
\]
Expanding \(\operatorname{Cov}_\nu(T_B,\Delta_B\mid Y)\) in the two coefficient coordinates and applying the third mixed-moment inequality gives
\[
|\operatorname{Cov}_\nu(T_B,\Delta_B\mid Y)|
\le2\mathfrak M_{3,n}T_{n,r}U_{n,r}X_n.
\tag{14}
\]
Since
\[
\operatorname{Var}_\nu(T_B+\Delta_B\mid Y)-\operatorname{Var}_\nu(T_B\mid Y)
=\operatorname{Var}_\nu(\Delta_B\mid Y)
+2\operatorname{Cov}_\nu(T_B,\Delta_B\mid Y),
\]
equations (9) and (14) prove the stated finite bound \(R_{n,r}^{\mathrm{ad}}\).

For the oracle identity, put \(q_B=B^{\top}\Omega\Pi^{-1}Y\).  Then
\[
q_B\in\operatorname{range}\{(\Omega^{1/2}B)^{\top}\}
=\operatorname{range}(S_B),
\]
and therefore \(S_BS_B^+q_B=q_B\).  Completing the quadratic variance expression gives
\[
\operatorname{Var}_\nu(\widehat\tau_B^{\mathrm{or},+}\mid Y)
=n^{-2}Y^{\top}\left\{\Pi^{-1}\Omega\Pi^{-1}
-\Pi^{-1}\Omega B S_B^+B^{\top}\Omega\Pi^{-1}\right\}Y.
\tag{15}
\]
Under ass:iid-potential-pairs, the signed outcome vector satisfies
\[
\mathbb E_{H^{\otimes n}}(YY^{\top})=Q(c_H).
\tag{16}
\]
Taking the trace expectation in (15), using (16), and invoking \(\texttt{def:completion-pseudoinverse-loss}\) proves exactly \(\ell_B^+(c_H)\).

We next justify the selected-basis assertion without inferring stability from independence.  Conditional on the pilot, \(\widetilde B\) is fixed.  On the event \(\{\operatorname{rank}(S_{\widetilde B})=r\}\), condition \((S)\) places its realized class in \(\mathcal S_{r,n}(\lambda_r,\overline a)\), so the fixed-basis calculation above applies with that \(r\).  On the rank-zero event, (2) makes the error exactly zero.  Assumption ass:pilot-main-outcome-independence permits the conditional design and outcome expectations used above; integrating the conditional inequalities over the pilot proves the displayed selected-basis bound.  For the encoded CAD selector, \(\texttt{thm:encoded-cad-pilot-selector-certificate}\) gives exact rank and spectral isolation, and exact algebraic comparison checks both \(\lambda_+(S_B)\ge n\lambda_r\) and the coordinate envelope.  Thus its record verifies \((S)\) when, and only when, those displayed checks pass.  No use of pilot independence supplies either check.

For the rate bounds, nonadjacent coordinates in a dependency graph are independent and hence uncorrelated.  Thus every nonzero covariance entry in column \(j\) lies in the closed dependency neighborhood of \(j\), proving \(s_n\le d_n\).  A centered triple moment vanishes when one distinct coordinate is separated from the other two; counting the remaining ordered triples gives at most \(8(2n)d_n^2\) possibilities.  A contributing ordered quadruple either has a connected dependency support or splits into two components of size two; the crude union count is at most \(64(2n)^2d_n^2\).  Since \(|Z_j|\le1\), expansion of the sign-weighted sums therefore gives
\[
\mathfrak M_{3,n}\le8(2n)d_n^2,
\qquad
\mathfrak M_{4,n}\le64(2n)^2d_n^2.
\tag{17}
\]
Substitution of \(d_n\le\overline d n^\alpha\) and \(\pi_{\min,n}\ge\underline\pi n^{-\beta}\) into the definitions of \(U_{n,r},X_n,T_{n,r}\) gives, uniformly over the two fixed positive margins,
\[
R_{n,r}^{\mathrm{ad}}
\le K_{\mathrm{ad}}
\{n^{4\alpha+4\beta-2}
+n^{(7/2)\alpha+3\beta-2}\}.
\tag{18}
\]
Because \((7/2)\alpha+3\beta\le4\alpha+4\beta\) and ass:zhao-first-order-rate states \(4\alpha+4\beta<1\), equation (18) gives \(\sup_r nR_{n,r}^{\mathrm{ad}}\to0\).  At rank two the positive eigenvalue is the smallest eigenvalue, so the definitions give \(\mathcal S_{2,n}(\kappa,\overline a)=\iota(\mathcal B_{\kappa,\overline a})\), and the ordinary-inverse lemma is recovered.

Finally, every displayed \(S(b_1,b_0)\) is positive semidefinite.  A positive-semidefinite \(2\times2\) matrix has rank two exactly when its determinant is positive, rank one exactly when its determinant is zero and trace positive, and rank zero exactly when its trace is zero.  Therefore a rank-two class exists exactly when the displayed determinant is positive for some pair.  A positive rank exists exactly when one diagonal block has a vector of positive quadratic form, because using that vector in one arm produces rank at least one.  If both diagonal blocks vanish, positive semidefiniteness of \(\Omega\) forces the cross block to vanish, so every control scale is zero; the converse is immediate.  These observations prove the four design-side and three individual-rank certificates. \(\square\)